% !TEX root = ../main.tex
\chapter{モノイドとモノイダルな見方：結合できる処理}
\label{chap:monoids-monoidal}

\begin{chapterabstract}
本章では、モノイドを「空の値」と「結合できる操作」を持つ構造として導入する。
文字列結合、ログ集約、リスト連結、設定マージ、数値集計など、多くの処理は同じ形で考えられる。
重要なのは、結合できることそのものではなく、結合の括弧を変えても意味が変わらないことと、空の入力を自然に扱えることである。
Lean 4 では、これらを等式として表し、結合律と単位元の法則が実務上どのような安全性を与えるかを確認する。
最後に、モノイダル圏を「逐次合成に加えて、対象と射をテンソルで並置する言語」として軽く眺める。
\end{chapterabstract}

\begin{learninggoals}
\item モノイドを「結合演算」と「単位元」を持つ構造として説明できる。
\item 結合律と単位元が、ログ集約や設定マージの安全性にどう関係するかを説明できる。
\item Lean 4 で、リスト結合や小さな設定マージの法則を等式として書ける。
\item モノイダル圏を、テンソルによる並置や資源合成を考えるための発展的な見方として位置づけられる。
\end{learninggoals}

\section{この章を読む理由}

ログを集める。設定をマージする。複数のバッチ結果を一つにまとめる。
このような処理では、「二つを一つに結合する」操作が繰り返し現れる。
最初は単純に見えるが、実務ではすぐに次のような問いが出てくる。

\begin{itemize}
\item 空のログを混ぜても結果は変わらないか。
\item バッチ結果を分割して集約しても、最後の結果は同じか。
\item 並列に処理した結果を後で結合しても安全か。
\item 設定マージの括弧を変えても、同じ設定になるか。
\end{itemize}

図\ref{fig:ch16-monoid-overview}は、単位元と二項演算だけでなく、左右の単位律と結合律までがモノイドのデータであることを示す。

\FigurePlaceholder{fig-ch16-monoid-overview.png}{0.90}{モノイドの構成要素}{fig:ch16-monoid-overview}

前章までで、圏、同型、積・余積、関手、自然変換を見てきた。
それらは主に「型の間の変換」や「構造を保った変換」を扱うための言葉だった。
本章のモノイドは、少し見え方が違う。
一つの型の中で、値どうしを安全にまとめるための言葉である。

\section{モノイドは「空」と「結合」を持つ構造}

\subsection{まずはログで考える}

ログをリストとして表すとする。
一つの処理がログを返し、別の処理もログを返す。
二つの処理をつなげたら、二つのログを連結したくなる。

ここで、ログを一つの型 \texttt{Log} として見る。
空ログを \texttt{emptyLog}、結合を \texttt{combineLog} と呼ぶことにする。
この二つだけなら、単なる便利関数である。
モノイドになるためには、法則が必要である。

\begin{definitionbox}
型 \(M\) に対して、次の三つがあるとき、\(M\) はモノイドとして扱える。
\begin{itemize}
\item 単位元 \(e : M\)。これは「空」や「何もない値」を表す。
\item 結合演算 \(\diamond : M \to M \to M\)。これは二つの値を一つにまとめる。
\item 単位元の法則と結合律。
\end{itemize}
法則は次の形で書ける。
\[
  e \diamond x = x, \qquad
  x \diamond e = x, \qquad
  (x \diamond y) \diamond z = x \diamond (y \diamond z).
\]
\end{definitionbox}

\subsection{モノイドは「便利な結合関数」ではない}

結合関数があるだけでは、まだモノイドとは呼べない。
たとえば、二つの設定ファイルをマージする関数を作ったとしても、その関数がどのような法則を満たすかを確認しなければならない。

\begin{warningbox}
「二つをまとめる関数がある」ことと、「モノイドとして安全に使える」ことは違う。
モノイドとして使いたいなら、少なくとも単位元と結合律を確認する必要がある。
特に設定マージは、競合解決の仕様によって法則が壊れることがある。
\end{warningbox}

本章では、すべての現実的な設定マージがモノイドになるとは主張しない。
むしろ、どの条件を満たせばモノイドとして扱えるかを明示するために、Lean で小さなモデルを作る。

\section{文字列結合、ログ結合、設定のマージ}

ソフトウェアでよく出会うモノイド的な構造を整理しておく。
次の表では、「値の型」「空の値」「結合演算」「注意点」を並べている。

\begin{small}
\begin{longtable}{@{}p{0.17\linewidth}p{0.18\linewidth}p{0.21\linewidth}p{0.32\linewidth}@{}}
\toprule
例 & 空の値 & 結合演算 & 注意点 \\
\midrule
文字列 & 空文字列 & 文字列結合 & 順序は通常入れ替えられない。 \\
リスト & 空リスト & リスト連結 & 要素順序を保つ。 \\
ログ & 空ログ & ログ連結 & 並列集約では時刻や順序の扱いを設計する。 \\
自然数の合計 & \(0\) & 加算 & 順序を入れ替えてもよい例。 \\
自然数の積 & \(1\) & 乗算 & 空の積を \(1\) と見られる。 \\
設定 & 空設定 & マージ & 競合解決が結合律を満たすか確認が必要。 \\
\bottomrule
\end{longtable}
\end{small}

各処理が部分結果を返し、それらを同じ演算で結合できれば、分割処理やツリー状の集約を設計しやすい。
ただし、分割後の結果をどの順序で並べるかは、結合律とは別に決める必要がある。

\subsection{順序を変えてよいとは限らない}

モノイドで保証されるのは、括弧の付け方を変えてよいことである。
つまり、次の二つは同じ結果になる。

\[
  (a \diamond b) \diamond c
  =
  a \diamond (b \diamond c)
\]

しかし、\(a \diamond b = b \diamond a\) は別の法則である。
これは可換性と呼ばれる。
モノイド一般には要求されない。

\begin{warningbox}
ログ連結では、\texttt{frontend ++ backend} と \texttt{backend ++ frontend} は通常同じではない。
モノイドだからといって、並列処理の結果を任意の順序で混ぜてよいわけではない。
順序を入れ替えたい場合は、可換性やソートなど、別の仕様が必要になる。
\end{warningbox}

\section{結合律と単位元}

\subsection{結合律は括弧を忘れさせる}

三つのログ \texttt{a}, \texttt{b}, \texttt{c} を考える。
まず \texttt{a} と \texttt{b} を結合してから \texttt{c} を結合する場合と、\texttt{b} と \texttt{c} を先に結合してから \texttt{a} を結合する場合がある。
結合律は、この二つが同じ結果になるという約束である。

図\ref{fig:ch16-associativity}の二つの木は、葉の順序 \(a,b,c\) を保ったまま、結合する組だけを変えている。
結合律は、この二つの評価結果が一致することを述べる。

\FigurePlaceholder{fig-ch16-associativity.png}{0.90}{結合律}{fig:ch16-associativity}

この法則があると、実装上の分割単位を変えやすくなる。
たとえば、次の二つの集約計画は同じ意味として扱える。

\begin{itemize}
\item \texttt{((batch1 ++ batch2) ++ batch3)}
\item \texttt{(batch1 ++ (batch2 ++ batch3))}
\end{itemize}

これは、バッチ処理の粒度やツリー集約の形を変えてもよい、という設計上の余裕を生む。
ただし、要素の順序そのものを変えているわけではない。

\subsection{単位元は空ケースを扱いやすくする}

集約処理でよく問題になるのが、入力が空のときである。
ログが一件もなければ空ログを返したい。
合計なら \(0\) を返したい。
積なら \(1\) を返したい。

単位元があると、空入力を例外扱いしなくてよい。
初期値として単位元を置き、あとは同じ結合演算を繰り返せばよい。

\section{Lean 4 で小さなモノイド法則を使う}

ここでは、Mathlib の高度な代数ライブラリには進まない。
まずは標準的なリスト連結を使い、ログ結合の法則を等式として確認する。

Lean では、\texttt{[]} と \texttt{++} について左右の単位律と結合律を定理として書き、すべてのログに成り立つことを確認する。

\begin{listing}[htbp]
\begin{minted}{lean4}
namespace Chapter16

abbrev Log := List String

def emptyLog : Log := []

def combineLog (a b : Log) : Log :=
  a ++ b

#eval combineLog ["start"] ["end"]  -- => ["start", "end"]

theorem combineLog_empty_left (xs : Log) :
    combineLog emptyLog xs = xs := by
  rfl

theorem combineLog_empty_right (xs : Log) :
    combineLog xs emptyLog = xs := by
  simp [combineLog, emptyLog]

theorem combineLog_assoc (a b c : Log) :
    combineLog (combineLog a b) c = combineLog a (combineLog b c) := by
  simp [combineLog]
\end{minted}
\caption{\texttt{combineLog} のモノイド則}
\label{lst:ch16_log_laws}
\end{listing}

このコードで重要なのは、\texttt{combineLog} の実装そのものではない。
本当に重要なのは、\texttt{combineLog\_assoc} が仕様として残ることである。
この定理があれば、ログ集約の括弧を変えるリファクタリングを、単なる気分ではなく等式として説明できる。

\subsection{自作の小さなモノイド構造}

次に、「空」「結合」「法則」をまとめた小さな構造体を作る。
これは Mathlib の本格的な \texttt{Monoid} を置き換えるものではない。
本章の目的に合わせた、学習用の最小モデルである。

\begin{listing}[htbp]
\begin{minted}{lean4}
structure SimpleMonoid (M : Type) where
  empty : M
  append : M -> M -> M
  empty_left : forall x, append empty x = x
  empty_right : forall x, append x empty = x
  append_assoc : forall x y z,
    append (append x y) z = append x (append y z)

-- List Nat は、空リストとリスト連結でモノイドになる。
def listNatMonoid : SimpleMonoid (List Nat) where
  empty := []
  append := fun xs ys => xs ++ ys
  empty_left := by
    intro xs
    rfl
  empty_right := by
    intro xs
    simp
  append_assoc := by
    intro xs ys zs
    simp
\end{minted}
\caption{\texttt{SimpleMonoid}}
\label{lst:ch16_simple_monoid}
\end{listing}

ここでは、\texttt{SimpleMonoid} のフィールドに証明を含めている。
この設計は、「この型には結合関数があります」という宣言ではなく、「この型は指定した法則を満たします」という宣言になる。

\section{サンプル：ログ集約・設定マージの順序安全性}

\subsection{ログ集約では、分割の仕方を変えたい}

ログ収集では、全ログを一度に集められるとは限らない。
サービス単位、時間帯単位、リクエスト単位で分割されることが多い。
後でそれらをまとめるとき、括弧の付け方に依存しないなら、集約計画を変更しやすくなる。

\begin{listing}[htbp]
\begin{minted}{lean4}
def planLeft (a b c : Log) : Log :=
  combineLog (combineLog a b) c

def planRight (a b c : Log) : Log :=
  combineLog a (combineLog b c)

theorem log_plan_safe (a b c : Log) :
    planLeft a b c = planRight a b c := by
  simpa [planLeft, planRight] using combineLog_assoc a b c
\end{minted}
\caption{\texttt{log\_plan\_safe}}
\label{lst:ch16_log_plan}
\end{listing}

ここで証明しているのは、ログの順序を入れ替えてよいという性質ではない。
証明しているのは、同じ順序のまま括弧の付け方だけを変えてよい、という性質である。
これが、モノイドの結合律が与える安全性である。

\subsection{設定マージは法則を確認してから使う}

設定マージもモノイド的に扱えることがある。
ただし、現実の設定マージは複雑になりやすい。
ここでは、説明のために小さな設定型を作る。
設定には、追加リトライ回数とタグ一覧があるとする。
マージは、リトライ回数を足し、タグを連結する操作とする。

\begin{listing}[htbp]
\begin{minted}{lean4}
structure Config where
  retry : Nat
  tags : List String
  deriving Repr

def emptyConfig : Config :=
  { retry := 0, tags := [] }

def mergeConfig (a b : Config) : Config :=
  { retry := a.retry + b.retry,
    tags := a.tags ++ b.tags }

#eval mergeConfig { retry := 1, tags := ["api"] }
                  { retry := 2, tags := ["db"] }
-- => { retry := 3, tags := ["api", "db"] }

theorem mergeConfig_empty_left (c : Config) :
    mergeConfig emptyConfig c = c := by
  cases c
  simp [mergeConfig, emptyConfig]

theorem mergeConfig_empty_right (c : Config) :
    mergeConfig c emptyConfig = c := by
  cases c
  simp [mergeConfig, emptyConfig]
\end{minted}
\caption{\texttt{Config} と \texttt{mergeConfig}}
\label{lst:ch16_config_merge}
\end{listing}

単位元の法則は、\texttt{emptyConfig} を前後どちらからマージしても元の設定が戻ることを述べている。
Lean では、構造体を \texttt{cases} で分解し、各フィールドの等式に落とす。
\texttt{retry} については \(0\) の加算、\texttt{tags} については空リストの連結が使われる。

\begin{listing}[htbp]
\begin{minted}{lean4}
theorem mergeConfig_assoc (a b c : Config) :
    mergeConfig (mergeConfig a b) c =
    mergeConfig a (mergeConfig b c) := by
  cases a
  cases b
  cases c
  simp [mergeConfig, Nat.add_assoc, List.append_assoc]

end Chapter16
\end{minted}
\caption{\texttt{mergeConfig\_assoc}}
\label{lst:ch16_config_assoc}
\end{listing}

この小さなモデルでは、設定マージは結合律を満たす。
そのため、\texttt{(base <> env) <> service} と \texttt{base <> (env <> service)} は同じ意味になる。

図\ref{fig:ch16-config-merge}では、空設定が左右の単位元となり、三つの設定の結合順だけを変えた二経路が同じ設定に到達する。
これはフィールドごとの加算とリスト連結が、それぞれ対応する法則を満たすためである。

\FigurePlaceholder{fig-ch16-config-merge.png}{0.90}{設定マージの法則}{fig:ch16-config-merge}

\subsection{このモデルが証明していないこと}

この証明は、本番システムのすべての設定マージが安全だと主張しているわけではない。
証明したのは、ここで定義した \texttt{Config} と \texttt{mergeConfig} に関する性質だけである。

\begin{warningbox}
本番の設定マージでは、キーの競合、優先順位、環境変数の展開、機密情報の扱い、外部ファイルの読み込み順序などが関わる。
Lean の小さなモデルを使うときは、「どの仕様だけをモデル化したか」を必ず明記する必要がある。
\end{warningbox}

\section{発展：モノイダル圏は「テンソルによる並置」の言語}

モノイダル圏は、圏 \(\mathcal C\) 上の二変数関手
\(\otimes : \mathcal C \times \mathcal C \to \mathcal C\)、単位対象、結合子、左右の単位子を持ち、それらが整合性条件を満たす構造である~\cite{riehl-category-theory-in-context}。
この定義だけから、実行時の並列性、独立性、交換可能性は導かれない。

ここまでは、一つの型の中で値を結合するモノイドを見てきた。
圏論では、さらに一般化して、対象や射を横に並べて合成する構造を考える。
これがモノイダル圏の入口である。

普通の圏では、たとえば \texttt{A -> B} と \texttt{B -> C} を逐次合成する。
モノイダル圏では、これに加えて二つの射をテンソルで並置できる。
型と関数を直積で組み合わせる場合なら、\texttt{A -> B} と \texttt{C -> D} から
\texttt{(A, C) -> (B, D)} を作る操作が具体例になる。

ソフトウェアでは、次のような場面をモデル化するときの構造上の見方として役に立つ。

\begin{itemize}
\item 二つの入力を別々に変換し、結果をペアとしてまとめる。
\item 独立性を別途確認した複数のサービス呼び出しの入出力を組にする。
\item 画像処理やETLで、別々の列やチャネルへの変換を一つの積上の関数として表す。
\item リソースや権限を合成し、全体として使える処理を作る。
\end{itemize}

図\ref{fig:ch16-monoidal-category}は、出力を次の入力へ渡す逐次合成と、二つの射をテンソルで組み合わせる操作を別の方向で描いている。
後者の横向き配置は、同時実行を表す記号ではない。

\FigurePlaceholder{fig-ch16-monoidal-category.png}{0.90}{逐次合成とテンソル}{fig:ch16-monoidal-category}

本章では、モノイダル圏の厳密な定義には進まない。
ここでは、圏の合成が「処理を縦につなぐ」言語だったのに対して、モノイダルな見方は「処理を横に並べる」言語を加える、と理解すればよい。
本格的には、対象どうしを並べるテンソル積、単位対象、結合の自然な同一視などを扱う。

\begin{warningbox}
「横に並べる」は図式上・構造上の表現であり、taskを同時実行するという意味ではない。実行時並列に読み替えるには、共有状態がないこと、失敗・取消・順序の観測が独立であることなど、追加の仕様が必要である。また一般のモノイダル圏は対称とは限らないため、左右を交換できるとも限らない。
\end{warningbox}

\section{よくある誤解}

\subsection{モノイドなら順序を入れ替えてよいのか}

よくない。
モノイドは結合律を要求するが、可換性は要求しない。
文字列、リスト、ログの結合は、典型的に順序を入れ替えられない。
順序を入れ替える処理を安全にしたいなら、可換モノイド、ソート済み表現、集合化、重複の扱いなど、追加の仕様が必要になる。

\subsection{単位元はただのデフォルト値か}

単なるデフォルト値ではない。
単位元は、結合演算と一緒に使ったときに値を変えない、という法則を満たす値である。
たとえば、空文字列は文字列結合の単位元だが、\texttt{"default"} は通常そうではない。

\subsection{設定マージは常にモノイドになるのか}

ならない。
設定マージの仕様による。
右優先、左優先、優先度つき、時刻つき、エラー蓄積つきなど、設計によって法則は変わる。
実務では、「モノイドとして扱えるはず」と決め打ちするのではなく、空の値と結合律を明示的に確認するのが安全である。

\section{本章のまとめ}

モノイドは、単位元と結合演算を持ち、単位元の法則と結合律を満たす構造である。

ログ、リスト、文字列、合計、設定マージなど、多くの実務上の集約処理はモノイド的に設計できる。

結合律は、括弧の付け方を変えてもよいことを保証するが、順序を入れ替えてよいことは保証しない。

単位元は、空入力や初期値を自然に扱うための仕様である。

Lean 4 では、モノイド法則を等式として書き、ログ集約や設定マージの小さな安全性を確認できる。

モノイダル圏は、逐次合成に加えてテンソルによる並置を扱う発展的な言語として理解できる。
